% Generated draft: requires mathematical and source review.
\documentclass[11pt]{article}
\usepackage[margin=25mm]{geometry}
\usepackage{fontspec}
\setmainfont{Latin Modern Roman}
\usepackage{amsmath,amssymb,mathtools,unicode-math}
\setmathfont{Latin Modern Math}
\usepackage{enumitem,needspace,xurl,hyperref}
\setlength{\parindent}{0pt}
\setlength{\parskip}{5pt}
\setlength{\emergencystretch}{3em}
\newcommand{\N}{\mathbb{N}}
\newcommand{\Z}{\mathbb{Z}}
\newcommand{\Q}{\mathbb{Q}}
\newcommand{\R}{\mathbb{R}}
\newcommand{\C}{\mathbb{C}}
\newcommand{\F}{\mathbb{F}}
\newcommand{\A}{\mathbb{A}}
\newcommand{\PP}{\mathbb{P}}
\newcommand{\E}{\mathbb{E}}
\newcommand{\eps}{\varepsilon}
\newcommand{\dd}{\,\mathrm{d}}
\begin{document}
\section*{UnsolvedMath ID: 20000460}
\textbf{Source problem code:} catalogue entry
\subsection{Definitions and mathematical statement}
For a degree-$d$ plane curve $C\subset\PP^2_{\Q}$ define
\[H([a:b:c])=\max\{|a|,|b|,|c|\}\quad\text{for primitive integer coordinates},\qquad N(C;B)=\#\{P\in C(\Q):H(P)\le B\}.\]
The repeated $d$ in $N(C;d)$ is intentional: the height cutoff equals the degree. The original does not explicitly impose irreducibility or geometric integrality.

\textbf{Original question (Heath-Brown).} Improve the uniform bound $N(C;d)\ll d^3$, given examples with $N(C;d)\gg d^2$. A dated remark by Roger Heath-Brown, September 25, 2004, says that $d^{5/2}$ is easy and that the challenge is to improve that exponent. Uniformity in the curve is essential; a bound with an arbitrary constant depending on $C$ would not express this degree-varying question.

\textbf{Attached unreviewed answer in the intended integral-curve reading.} The aid attributes to Binyamini, Cluckers and Kato (2025) a uniform estimate
\[N(C;d)\le c\,d^2d^{2/d}(\log d)^{12}=O\bigl(d^2(\log d)^{12}\bigr)=o(d^{5/2}),\qquad d\ge2,\]
for the intended integral class of curves, with an absolute constant $c$. Here ``integral'' concerns irreducibility/reducedness, not a count only of affine integer points. The exact theorem's irreducible versus geometrically integral hypotheses and logarithm conventions are not supplied; the case $d=1$ cannot literally use the displayed $\log d$ expression and must be treated separately.

The aid additionally claims that for every integer $d\ge1$ one can choose an integer $\lambda_d$ with $1\le\lambda_d\le(d-1)^2+1$ such that
\[\prod_{i=1}^d(X-iZ)+\lambda_d\prod_{j=1}^d(Y-jZ)=0\]
is smooth and geometrically integral over $\Q$. It contains every grid point $[i:j:1]$ with $1\le i,j\le d$, all of height at most $d$, and hence has at least $d^2$ counted points. Thus the quadratic scale is reported already in a smooth geometrically integral family.

In contrast, the aid claims that for every sufficiently large $d$ there is a reduced reducible degree-$d$ curve with $N(C;d)=\Omega(d^{8/3})$, using a coefficient-height and line-lattice construction not supplied in the summary. This shows why the omitted irreducibility convention matters: that reducible growth exceeds even the dated $d^{5/2}$ bound. The theorem attribution and both construction claims require independent review.
\subsection{Short English statement}
The degree and height bound are both $d$. Heath-Brown's dated challenge is to beat a uniform $d^{5/2}$ bound. The aid reports a 2025 quadratic bound up to logarithms for integral curves, a smooth $d^2$ grid construction, and faster reducible examples showing that the omitted irreducibility assumption is essential.
\subsection{Sources}
\begin{itemize}
\item Local UnsolvedMath registry, permanent ID 20000460, source code catalogue entry; \texttt{aim-algebraic-number-theory-notes.json}, notes[111], tag \texttt{question}.
\item AIM workshop \emph{Rational and integral points on higher dimensional varieties}, Question 27 (Heath-Brown): \url{https://aimath.org/WWN/qptsurface2/qptsurface2.pdf}. Remark by Roger Heath-Brown dated September 25, 2004. The supplied verification restores $\ll$, $\gg$ and exponents $3,2,5/2$ and confirms the repeated $d$. No PDF was consulted.
\item Attached research attempt 2 attributes the known solution to Binyamini--Cluckers--Kato (2025) without a title, identifier or theorem number. Result \texttt{partial\_result}, problem status at run \texttt{solved}, rigor label \texttt{proved}, novelty confidence low. It is machine generated, not peer reviewed, and requires independent verification before citation. Original literature field empty.
\end{itemize}
\end{document}
